% introduction.tex
\section{Introduction}
\label{sec:intro}

Factor models are the workhorse of asset pricing, risk management, and portfolio construction. Practitioners and academics rely on a small number of common factors---market, size, value, momentum---to explain the cross-section of returns, allocate risk budgets, and evaluate investment performance. Yet the dominant factor models are built from equity sorts and evaluated almost exclusively on stock returns. Bonds, REITs, and emerging market assets collectively represent trillions of dollars in assets under management, and for these asset classes the standard Fama-French factors explain as little as 11\% of return variation out of sample.% TODO: verify

Deep learning has transformed factor modeling in recent years. Conditional autoencoders learn nonlinear mappings from asset characteristics to time-varying factor loadings \citep{GuKellyXiu2021autoencoder}, while GAN-based stochastic discount factor estimation \citep{ChenPelgerZhu2024deep}, risk-premium PCA \citep{LettauPelger2020factors}, and instrumented PCA \citep{KellyPruittSu2019ipca} push predictive performance further on US equities. VAE-based extensions introduce probabilistic factor structure \citep{Duan2022factorvae}, and attention mechanisms capture cross-asset interactions \citep{Epstein2025attention}. These models achieve impressive out-of-sample $R^2$ and Sharpe ratios on individual stock returns.

Two critical gaps remain, however. First, most deep factor models produce latent factors that are abstract statistical objects---they cannot be directly traded or economically interpreted. While \citet{GuKellyXiu2021autoencoder} originally proposed a no-arbitrage weight structure, subsequent extensions relax this constraint in favor of richer latent representations, severing the link between learned factors and implementable trading strategies. Second, virtually all deep factor models are trained and evaluated exclusively on US equities. Whether these architectures generalize to bonds, REITs, and emerging market assets---where the underlying return drivers differ fundamentally from equity risk premia---is unknown. Meanwhile, fund-level alternative data such as portfolio holdings and sector allocations from regulatory filings remains unexploited as conditioning information for factor construction.

We propose the Autoencoding Factor Model (AFM), which addresses both gaps simultaneously. The key architectural innovation is a \emph{portfolio bottleneck}: rather than learning free latent codes, the AFM constrains each factor to be the realized return of an explicit portfolio whose weights are produced by a neural network and mapped to valid allocations via softmax (long-only) or a softmax-minus-softmax construction (market-neutral). Factor loadings and intercepts are generated by separate networks conditioned on fund-level characteristics---economic sector holdings and asset class allocations derived from SEC N-PORT filings---together with a learned market-state embedding that compresses a 252-day return window. Because all portfolio weights and loadings depend only on information available at time $t-1$, every factor is tradable without look-ahead bias. The entire system is trained end-to-end via return reconstruction loss, inheriting the scalability of modern deep learning while producing factors that connect directly to the classical tradition of mimicking portfolios \citep{Huberman1987mimicking} and parametric portfolio policies \citep{Brandt2009parametric}.

\paragraph{Contributions.}
\begin{itemize}
    \item \textbf{Portfolio bottleneck architecture.} We propose a neural autoencoder where each latent factor is constrained to be the return of a tradable portfolio with explicit, no-look-ahead weights, bridging deep latent factor models \citep{GuKellyXiu2021autoencoder} and classical portfolio-based factor construction \citep{Huberman1987mimicking}. This is, to our knowledge, the first autoencoder factor model that enforces tradability as a hard architectural constraint rather than a soft regularizer.
    \item \textbf{Conditional factor structure via fund-level alternative data.} Portfolio weights and factor loadings are conditioned on mutual fund characteristics (economic sector holdings, asset class allocations) derived from regulatory filings, enabling interpretable, time-varying factor exposures that reflect the economic composition of each fund.
    \item \textbf{Cross-asset factor learning.} Trained on approximately 1{,}000 US mutual funds across 10 asset classes---spanning large, mid, and small cap equities, international and emerging market equities, US aggregate and high-yield bonds, emerging market debt, municipal bonds, and REITs---the AFM achieves 85.0\% average out-of-sample $R^2$ compared to 68.4\% for the Fama-French four-factor model \citep{FamaFrench2015fivefactor, Carhart1997persistence}, with the largest gains on non-equity assets where traditional equity factors have near-zero explanatory power.% TODO: verify
    \item \textbf{Comprehensive ablation study.} We systematically evaluate architecture choices (MLP vs.\ residual MLP), training universes (single-asset-class vs.\ multi-asset), and their interactions. We find that simpler architectures and broader training universes yield superior generalization, with multi-asset training improving average $R^2$ by 12 percentage points over large-cap-only training.% TODO: verify
\end{itemize}

\paragraph{Paper Organization.}
The remainder of this paper is organized as follows. Section~\ref{sec:related} reviews related work on deep factor models, autoencoder extensions, and fund-level machine learning. Section~\ref{sec:method} formalizes the AFM architecture and training procedure. Section~\ref{sec:experiments} presents our main results and ablation studies. Section~\ref{sec:conclusion} concludes.
